\chapter{Analyse et spécification des besoins}
\label{chap:analyse}

\sectionsansnum{Introduction}

Le chapitre précédent a établi ce qu'il fallait construire et pourquoi. Celui-ci
formalise la commande. Nous identifions d'abord les acteurs du système et leurs
responsabilités, puis nous exprimons les besoins fonctionnels et non fonctionnels. Ces
besoins sont ensuite traduits en un \emph{Product Backlog} de soixante récits
utilisateur, priorisés et estimés. Nous présentons enfin le modèle de données,
l'environnement de travail, l'architecture globale et la répartition du backlog en six
sprints, qui structure les six chapitres suivants.

\section{Identification des acteurs}

Un acteur désigne toute entité qui interagit avec le système sans en faire partie. On
distingue les acteurs \emph{principaux}, qui déclenchent les cas d'utilisation pour en
retirer une valeur, des acteurs \emph{secondaires}, sollicités par le système pour
mener une opération à bien. Le tableau~\ref{tab:acteurs} recense les quatre acteurs de
\hypnoria{}.

\begin{table}[H]
\centering

\small
\begin{tabular}{@{}L{3.0cm}cL{8.2cm}@{}}
\toprule
\textbf{Acteur} & \textbf{Type} & \textbf{Responsabilités} \\
\midrule
Médecin du sommeil & principal &
Constituer et suivre son registre de patients, lancer l'analyse d'un enregistrement,
interpréter l'hypnogramme et la prédiction de sévérité, annoter, produire le compte
rendu, solliciter un confrère \\
\addlinespace[3pt]
Administrateur & principal &
Créer et administrer les comptes médecins, gérer les échéances de licence, rattacher
les praticiens à un établissement, superviser l'activité et consulter le journal
d'audit \\
\addlinespace[3pt]
Moteur d'inférence & secondaire &
Charger les modèles, prétraiter le signal, produire la séquence de stades, calculer
les métriques et les attributions d'explication \\
\addlinespace[3pt]
Service de stockage & secondaire &
Recevoir les fichiers volumineux, les conserver de façon durable et délivrer des liens
d'accès temporaires \\
\bottomrule
\end{tabular}
\caption{Acteurs du système et leurs responsabilités}
\label{tab:acteurs}
\end{table}

Deux remarques sur ce découpage. Le \textbf{moteur d'inférence} est modélisé comme un
acteur distinct bien qu'il réside dans le même processus que l'interface applicative :
ce choix matérialise le fait qu'il pourrait être déporté sur une machine dédiée sans
modifier les cas d'utilisation. Le \textbf{service de stockage}, lui, est réellement
externe.

\section{Étude des besoins}

\subsection{Les besoins fonctionnels}

Les besoins fonctionnels ont été organisés en neuf \emph{épics}, regroupements
thématiques de fonctionnalités partageant un même objectif métier. Cette structure,
présentée au tableau~\ref{tab:epics}, sert de colonne vertébrale au backlog.

\begin{table}[H]
\centering

\small
\begin{tabular}{@{}clL{7.4cm}cc@{}}
\toprule
\textbf{Épic} & \textbf{Intitulé} & \textbf{Portée} & \textbf{Récits} &
\textbf{Points} \\
\midrule
E1 & Authentification & Connexion, jetons, hachage, contrôle des rôles, application des
licences & 5 & 18 \\
\addlinespace[2pt]
E2 & Administration & Comptes médecins, cycle de vie, établissements, audit,
statistiques & 10 & 45 \\
\addlinespace[2pt]
E3 & Dossier patient & Registre des patients, examens, historique, tableau de bord
& 6 & 24 \\
\addlinespace[2pt]
E4 & Classification & Lecture EDF, prétraitement, modèles de base, métriques AASM
& 12 & 95 \\
\addlinespace[2pt]
E5 & Sévérité du SAOS & Empilement, extraction de marqueurs, prédiction, explications
& 8 & 65 \\
\addlinespace[2pt]
E6 & Visualisation & Hypnogramme, oxymétrie, annotation, normes de référence & 5 & 28 \\
\addlinespace[2pt]
E7 & Collaboration & Discussions rattachées à un examen, boîte de réception & 3 & 18 \\
\addlinespace[2pt]
E8 & Archivage et export & Stockage cloud, liens signés, compte rendu, exports
& 6 & 36 \\
\addlinespace[2pt]
E9 & Industrialisation & Conteneurisation, intégration continue, multilinguisme
& 5 & 29 \\
\midrule
\multicolumn{3}{@{}r}{\textbf{Total}} & \textbf{60} & \textbf{358} \\
\bottomrule
\end{tabular}
\caption{Regroupement des besoins fonctionnels en épics}
\label{tab:epics}
\end{table}

\subsection{Les besoins non fonctionnels}

Les exigences non fonctionnelles conditionnent l'acceptabilité du système autant que
ses fonctions. Elles sont d'autant plus contraignantes ici que le domaine --- la donnée
de santé --- impose des obligations propres. Le tableau~\ref{tab:bnf} les énonce avec,
pour chacune, le moyen de vérification retenu : un besoin non fonctionnel qu'on ne sait
pas vérifier n'est qu'une intention.

\begin{table}[H]
\centering

\small
\begin{tabular}{@{}L{2.7cm}L{5.6cm}L{5.6cm}@{}}
\toprule
\textbf{Critère} & \textbf{Exigence} & \textbf{Vérification} \\
\midrule
Sécurité &
Aucun accès sans jeton valide ; mots de passe jamais stockés en clair ; cloisonnement
strict entre rôles &
Revue des dépendances de sécurité sur chaque route ; tentative d'accès croisé entre
comptes \\
\addlinespace[3pt]
Confidentialité &
Aucun fichier patient accessible par une URL publique ou devinable &
Vérification que tout lien délivré est signé et expire \\
\addlinespace[3pt]
Traçabilité &
Toute action sensible horodatée avec son auteur, journal exportable &
Inspection du journal après un scénario complet \\
\addlinespace[3pt]
Performance &
Analyse d'une nuit complète en quelques minutes ; interface non bloquée pendant les
transferts &
Chronométrage sur enregistrements réels ; transfert en arrière-plan observable \\
\addlinespace[3pt]
Interopérabilité &
Accepter les fichiers EDF quel que soit le constructeur ; importer et exporter en
formats ouverts &
Résolution des canaux sur des enregistrements de provenances différentes \\
\addlinespace[3pt]
Ergonomie &
Parcours d'analyse guidé, sans formation préalable ; retour visuel sur les traitements
longs &
Passage d'un utilisateur non initié sur le parcours complet \\
\addlinespace[3pt]
Multilinguisme &
Interface intégralement disponible en français et en anglais &
Revue de l'absence de chaînes non traduites \\
\addlinespace[3pt]
Portabilité &
Déploiement complet sur une machine vierge sans configuration manuelle &
Démarrage de la pile sur un poste neuf en une commande \\
\addlinespace[3pt]
Maintenabilité &
Séparation nette des responsabilités ; ajout d'un modèle sans toucher au reste &
Ajout d'une variante de modèle à titre d'essai \\
\bottomrule
\end{tabular}
\caption{Besoins non fonctionnels et moyens de vérification}
\label{tab:bnf}
\end{table}

\section{Diagramme de cas d'utilisation global}

La figure~\ref{fig:uc-global} présente la vue d'ensemble des interactions entre les
acteurs et le système.

\begin{figure}[H]
\centering
\begin{tikzpicture}[
    x=1cm,y=1cm,
    uc/.style      = {ellipse, draw=hypnoblue!60, fill=hypnolight,
                      align=center, font=\scriptsize, text width=2.45cm,
                      inner sep=1.5pt, minimum height=0.92cm},
    lien/.style    = {hypnogray!75, line width=0.5pt},
    nomact/.style  = {font=\scriptsize\bfseries, align=center, text width=2.2cm},
  ]

  % ── Pictogramme d'acteur ────────────────────────────────────────────────
  \tikzset{acteur/.pic={
      \draw[hypnoblue,line width=0.8pt] (0,0) circle (0.17);
      \draw[hypnoblue,line width=0.8pt] (0,-0.17) -- (0,-0.66);
      \draw[hypnoblue,line width=0.8pt] (-0.28,-0.34) -- (0.28,-0.34);
      \draw[hypnoblue,line width=0.8pt] (0,-0.66) -- (-0.24,-1.04);
      \draw[hypnoblue,line width=0.8pt] (0,-0.66) -- (0.24,-1.04);}}

  % ── Frontière du système ────────────────────────────────────────────────
  \draw[hypnoblue!55,line width=0.9pt,rounded corners=3pt]
       (3.45,-6.40) rectangle (11.05,2.55);
  \node[font=\small\bfseries,text=hypnoblue,anchor=north] at (7.25,2.42)
       {Système \hypnoria{}};

  % ── Cas d'utilisation ───────────────────────────────────────────────────
  \node[uc] (auth)  at (7.25, 1.20) {S'authentifier};

  \node[uc] (pat)   at (5.35,-0.40) {Gérer ses patients};
  \node[uc] (ana)   at (5.35,-1.70) {Analyser un\\enregistrement};
  \node[uc] (osa)   at (5.35,-3.00) {Prédire la\\sévérité du SAOS};
  \node[uc] (annot) at (5.35,-4.30) {Annoter et\\exporter};
  \node[uc] (colab) at (5.35,-5.60) {Solliciter un\\confrère};

  \node[uc] (comp)  at (9.15,-0.40) {Gérer les comptes\\et licences};
  \node[uc] (etab)  at (9.15,-1.70) {Gérer les\\établissements};
  \node[uc] (audit) at (9.15,-3.00) {Consulter le\\journal d'audit};
  \node[uc] (stats) at (9.15,-4.30) {Suivre l'activité\\de la plateforme};

  % ── Acteurs principaux : libellé au-dessus, liaisons vers la droite ─────
  \pic at (1.15,0.00) {acteur};
  \node[nomact,anchor=south] at (1.15,0.42) {Médecin\\du sommeil};

  \pic at (13.35,0.00) {acteur};
  \node[nomact,anchor=south] at (13.35,0.42) {Administrateur};

  \foreach \c in {auth,pat,ana,osa,annot,colab}
    \draw[lien] (1.50,-0.20) -- (\c.west);

  \foreach \c in {auth,comp,etab,audit,stats}
    \draw[lien] (13.00,-0.20) -- (\c.east);

  % ── Acteurs secondaires : libellé en dessous ────────────────────────────
  \pic at (7.25,-7.55) {acteur};
  \node[nomact,anchor=north] at (7.25,-8.72) {Moteur\\d'inférence};

  \pic at (10.35,-7.55) {acteur};
  \node[nomact,anchor=north] at (10.35,-8.72) {Service de\\stockage};

  \draw[lien] (7.25,-7.28) -- (ana.east);
  \draw[lien] (7.25,-7.28) -- (osa.east);
  \draw[lien] (10.35,-7.28) -- (annot.east);

\end{tikzpicture}
\vspace{0.2cm}
\caption[Diagramme de cas d'utilisation global]%
        {Vue d'ensemble des cas d'utilisation de \hypnoria{}.}
\label{fig:uc-global}
\end{figure}

Le cas « Analyser un enregistrement » constitue le c\oe ur fonctionnel du système. La
figure~\ref{fig:uc-analyse} le décompose et fait apparaître les relations d'inclusion
--- sous-cas systématiquement exécutés --- et d'extension --- comportements optionnels.

\begin{figure}[H]
\centering
\begin{tikzpicture}[
    x=1cm,y=1cm,
    uc/.style      = {ellipse, draw=hypnoblue!60, fill=hypnolight,
                      align=center, font=\scriptsize, text width=2.5cm,
                      inner sep=1.5pt, minimum height=0.92cm},
    ucp/.style     = {ellipse, draw=hypnored!75, fill=hypnored!8,
                      align=center, font=\scriptsize\bfseries, text width=2.5cm,
                      inner sep=1.5pt, minimum height=0.92cm},
    lien/.style    = {hypnogray!75, line width=0.5pt},
    incl/.style    = {-{Latex[length=2mm]}, hypnoblue!75, line width=0.5pt, dashed},
    ext/.style     = {-{Latex[length=2mm]}, hypnored!70, line width=0.5pt, dashed},
    et/.style      = {font=\tiny\itshape, fill=white, inner sep=1.2pt},
    nomact/.style  = {font=\scriptsize\bfseries, align=center, text width=1.8cm},
  ]

  \tikzset{acteur/.pic={
      \draw[hypnoblue,line width=0.8pt] (0,0) circle (0.17);
      \draw[hypnoblue,line width=0.8pt] (0,-0.17) -- (0,-0.66);
      \draw[hypnoblue,line width=0.8pt] (-0.28,-0.34) -- (0.28,-0.34);
      \draw[hypnoblue,line width=0.8pt] (0,-0.66) -- (-0.24,-1.04);
      \draw[hypnoblue,line width=0.8pt] (0,-0.66) -- (0.24,-1.04);}}

  % ── Frontière ───────────────────────────────────────────────────────────
  \draw[hypnoblue!55,line width=0.9pt,rounded corners=3pt]
       (2.75,-4.55) rectangle (12.60,3.15);
  \node[font=\small\bfseries,text=hypnoblue,anchor=north] at (7.65,3.02)
       {Module d'analyse};

  % ── Cas central et cas déclenchés par le médecin ────────────────────────
  \node[ucp] (ana) at (4.75, 0.75) {Analyser un\\enregistrement};
  \node[uc]  (exp) at (4.75,-2.20) {Exporter les\\marqueurs};
  \node[uc]  (pdf) at (4.75,-3.65) {Produire le\\compte rendu};

  % ── Sous-cas inclus ─────────────────────────────────────────────────────
  \node[uc] (chn) at (10.30, 2.10) {Choisir le montage\\et la granularité};
  \node[uc] (pre) at (10.30, 0.85) {Prétraiter\\le signal};
  \node[uc] (cls) at (10.30,-0.40) {Classer les époques};
  \node[uc] (met) at (10.30,-1.75) {Calculer les\\métriques AASM};

  % ── Cas étendu ──────────────────────────────────────────────────────────
  \node[uc] (osa) at (10.30,-3.65) {Prédire la\\sévérité};

  % ── Acteur ──────────────────────────────────────────────────────────────
  \pic at (1.05,0.55) {acteur};
  \node[nomact,anchor=south] at (1.05,0.97) {Médecin};

  \draw[lien] (1.40,0.35) -- (ana.west);
  \draw[lien] (1.40,0.35) -- (exp.west);
  \draw[lien] (1.40,0.35) -- (pdf.west);

  % ── Relations d'inclusion (libellés étagés le long du faisceau) ──────────
  \draw[incl] (ana) -- node[et,text=hypnoblue!85,pos=0.30] {$\ll$include$\gg$} (chn);
  \draw[incl] (ana) -- node[et,text=hypnoblue!85,pos=0.44] {$\ll$include$\gg$} (pre);
  \draw[incl] (ana) -- node[et,text=hypnoblue!85,pos=0.58] {$\ll$include$\gg$} (cls);
  \draw[incl] (ana) -- node[et,text=hypnoblue!85,pos=0.72] {$\ll$include$\gg$} (met);

  % ── Relations d'extension ───────────────────────────────────────────────
  \draw[ext] (osa) -- node[et,text=hypnored!85,pos=0.5] {$\ll$extend$\gg$} (met);
  \draw[ext] (pdf) -- node[et,text=hypnored!85,pos=0.5] {$\ll$extend$\gg$} (exp);

\end{tikzpicture}
\vspace{0.2cm}
\caption[Raffinement du cas « Analyser un enregistrement »]%
        {Décomposition du cas d'utilisation central en sous-cas inclus et étendus.}
\label{fig:uc-analyse}
\end{figure}

Les tableaux~\ref{tab:desc-analyse} et~\ref{tab:desc-osa} décrivent textuellement les
deux cas les plus critiques. Un diagramme montre en effet la structure des interactions,
mais pas leur déroulement.

\begin{table}[H]
\centering

\small
\begin{tabular}{@{}L{3.0cm}L{11.4cm}@{}}
\toprule
\textbf{Identifiant} & UC-03 \\
\textbf{Acteur principal} & Médecin du sommeil \\
\textbf{Acteurs secondaires} & Moteur d'inférence, service de stockage \\
\textbf{Préconditions} & Le médecin est authentifié ; sa licence est valide ; il
dispose d'un fichier d'enregistrement au format EDF \\
\textbf{Postconditions} & Un examen est créé dans le dossier du patient, la séquence
de stades et les métriques sont enregistrées, le fichier brut est archivé \\
\midrule
\textbf{Scénario nominal} &
\begin{minipage}[t]{11.2cm}
\begin{enumerate}[leftmargin=1.1em,itemsep=1pt,topsep=2pt]
  \item Le médecin sélectionne le patient concerné.
  \item Il choisit le montage disponible, à deux ou cinq dérivations.
  \item Il choisit la granularité de classification, à trois ou cinq stades.
  \item Il dépose le fichier d'enregistrement.
  \item Le système décode le fichier, résout les dérivations et prétraite le signal.
  \item Le moteur d'inférence classe chaque époque et le système calcule les métriques.
  \item Le système affiche l'hypnogramme, les métriques et les alertes cliniques.
  \item Le système crée l'examen et lance l'archivage du fichier en arrière-plan.
\end{enumerate}
\end{minipage} \\
\midrule
\textbf{Scénarios alternatifs} &
\begin{minipage}[t]{11.2cm}
\begin{itemize}[leftmargin=1.1em,itemsep=1pt,topsep=2pt]
  \item \textbf{4a} --- Le fichier n'est pas au format attendu : le système refuse le
        dépôt et le signale.
  \item \textbf{5a} --- Une dérivation est absente : le système applique son repli
        documenté et poursuit.
  \item \textbf{5b} --- L'enregistrement est plus court qu'une époque : le traitement
        s'interrompt avec un message explicite.
  \item \textbf{8a} --- L'archivage échoue : les résultats restent consultables et le
        transfert est signalé en erreur.
\end{itemize}
\end{minipage} \\
\bottomrule
\end{tabular}
\caption{Description textuelle du cas « Analyser un enregistrement »}
\label{tab:desc-analyse}
\end{table}

\begin{table}[H]
\centering

\small
\begin{tabular}{@{}L{3.0cm}L{11.4cm}@{}}
\toprule
\textbf{Identifiant} & UC-04 \\
\textbf{Acteur principal} & Médecin du sommeil \\
\textbf{Acteur secondaire} & Moteur d'inférence \\
\textbf{Préconditions} & Une séquence de stades est disponible, issue d'une analyse ou
d'une saisie manuelle \\
\textbf{Postconditions} & Une classe de sévérité, ses probabilités et la décomposition
des facteurs sont associées à l'examen \\
\midrule
\textbf{Scénario nominal} &
\begin{minipage}[t]{11.2cm}
\begin{enumerate}[leftmargin=1.1em,itemsep=1pt,topsep=2pt]
  \item Le système extrait les marqueurs d'architecture depuis l'hypnogramme.
  \item Le médecin complète les données démographiques et oxymétriques.
  \item Il déclenche l'évaluation du risque.
  \item Le système construit le vecteur de variables et comble les valeurs absentes.
  \item Le moteur produit la classe de sévérité et ses probabilités.
  \item Le système calcule les contributions de chaque variable.
  \item Le système affiche la sévérité, les facteurs aggravants et protecteurs, et
        l'interprétation clinique.
\end{enumerate}
\end{minipage} \\
\midrule
\textbf{Scénarios alternatifs} &
\begin{minipage}[t]{11.2cm}
\begin{itemize}[leftmargin=1.1em,itemsep=1pt,topsep=2pt]
  \item \textbf{2a} --- Le médecin importe un rapport externe : le système en extrait
        les valeurs et préremplit le formulaire.
  \item \textbf{4a} --- Des variables cliniques manquent : le système applique la
        valeur médiane de la population d'apprentissage et le signale.
  \item \textbf{5a} --- Les modèles ne sont pas chargés : l'opération est refusée avec
        un message explicite.
\end{itemize}
\end{minipage} \\
\bottomrule
\end{tabular}
\caption{Description textuelle du cas « Prédire la sévérité du SAOS »}
\label{tab:desc-osa}
\end{table}

\section{Backlog du produit}

Le \emph{Product Backlog} est la liste ordonnée de tout ce que le produit doit offrir.
Chaque élément prend la forme d'un récit utilisateur suivant la formulation canonique :
\emph{en tant que} \dots, \emph{je veux} \dots, \emph{afin de} \dots. Cette
formulation force à énoncer le bénéficiaire et la finalité, et non seulement la
fonction.

Deux attributs complètent chaque récit. La \textbf{priorité} suit la méthode MoSCoW ---
\emph{Must have} pour ce sans quoi le produit n'a pas de sens, \emph{Should have} pour
ce qui lui donne sa valeur, \emph{Could have} pour ce qui l'améliore sans le
conditionner. L'\textbf{estimation} s'exprime en points de complexité sur la suite de
Fibonacci, échelle relative qui compare les récits entre eux plutôt que de prétendre
prédire une durée.

\begin{longtable}{@{}p{1.25cm}L{9.1cm}cp{0.75cm}p{0.9cm}@{}}
\caption{Product Backlog --- 60 récits utilisateur, 358 points}
\label{tab:backlog}\\
\toprule
\textbf{ID} & \textbf{Récit utilisateur} & \textbf{Prio.} & \textbf{Pts} &
\textbf{Spr.} \\
\midrule
\endfirsthead

\multicolumn{5}{@{}l}{\small\itshape Tableau~\ref{tab:backlog} --- suite}\\
\toprule
\textbf{ID} & \textbf{Récit utilisateur} & \textbf{Prio.} & \textbf{Pts} &
\textbf{Spr.} \\
\midrule
\endhead

\midrule
\multicolumn{5}{@{}r}{\small\itshape suite page suivante}\\
\endfoot

\bottomrule
\endlastfoot

\small
US-01 & Me connecter avec mon adresse et mon mot de passe pour accéder à la plateforme
& M & 5 & S4 \\
US-02 & Consulter mon profil pour vérifier mon rôle et mon rattachement & M & 2 & S4 \\
US-03 & Stocker les mots de passe hachés pour qu'aucun ne soit lisible en base
& M & 3 & S4 \\
US-04 & Vérifier le rôle sur chaque route pour cloisonner médecins et administrateurs
& M & 5 & S4 \\
US-05 & Refuser la connexion d'un compte suspendu ou dont la licence a expiré
& M & 3 & S5 \\

\epictitre{E2 --- Administration de la plateforme}
US-06 & Créer un compte médecin pour ouvrir l'accès à un nouveau praticien & M & 5 & S5 \\
US-07 & Lister les médecins avec leur statut et leur échéance de licence & M & 3 & S5 \\
US-08 & Suspendre ou réactiver un compte et fixer sa date d'expiration & M & 5 & S5 \\
US-09 & Réinitialiser le mot de passe d'un médecin pour débloquer un accès perdu
& S & 3 & S5 \\
US-10 & Créer des établissements et y rattacher des médecins & S & 5 & S5 \\
US-11 & Journaliser chaque action sensible avec son auteur et son horodatage
& M & 5 & S5 \\
US-12 & Consulter le journal d'audit pour enquêter sur un incident & M & 5 & S5 \\
US-13 & Exporter le journal d'audit pour un contrôle de conformité & S & 3 & S5 \\
US-14 & Consulter les statistiques globales pour suivre l'adoption & S & 3 & S5 \\
US-15 & Visualiser la chaîne de traitement des modèles & C & 8 & S6 \\

\epictitre{E3 --- Dossier patient et examens}
US-16 & Enregistrer un patient avec son âge, son IMC et son sexe & M & 3 & S5 \\
US-17 & Lister mes patients avec leur historique d'examens & M & 5 & S5 \\
US-18 & Consulter la fiche détaillée d'un patient et tous ses examens & M & 3 & S5 \\
US-19 & Créer un examen rattaché à un patient pour y associer les résultats
& M & 5 & S5 \\
US-20 & Mettre à jour la sévérité et les résultats d'un examen & M & 3 & S5 \\
US-21 & Voir la répartition des risques d'apnée de ma patientèle & S & 5 & S6 \\

\epictitre{E4 --- Classification automatique des stades}
US-22 & Lire un fichier EDF et en extraire les dérivations physiologiques & M & 8 & S1 \\
US-23 & Résoudre les canaux malgré l'hétérogénéité des conventions de nommage
& M & 5 & S1 \\
US-24 & Rééchantillonner, segmenter en époques de 30 s et normaliser & M & 8 & S1 \\
US-25 & Valider la chaîne sur un jeu public en attendant l'accès aux cohortes
& M & 8 & S1 \\
US-26 & Prétraiter les enregistrements de la cohorte clinique à l'échelle & M & 8 & S2 \\
US-27 & Classer chaque époque avec un réseau convolutif multi-échelle & M & 13 & S2 \\
US-28 & Classer chaque époque avec un réseau récurrent à pooling attentionnel
& M & 13 & S2 \\
US-29 & Classer chaque époque avec un Transformer par patchs & M & 13 & S2 \\
US-30 & Choisir la configuration d'analyse selon le montage disponible & S & 5 & S4 \\
US-31 & Inspecter les canaux d'un fichier avant de lancer l'analyse & S & 3 & S4 \\
US-32 & Calculer les métriques AASM de l'enregistrement & M & 8 & S4 \\
US-33 & Conserver les modèles en mémoire entre deux requêtes & S & 3 & S4 \\

\epictitre{E5 --- Ensemble et prédiction de la sévérité}
US-34 & Combiner les trois modèles par un méta-apprenant & M & 13 & S3 \\
US-35 & Extraire les marqueurs d'architecture du sommeil depuis l'hypnogramme
& M & 8 & S3 \\
US-36 & Prédire la sévérité du SAOS en quatre classes & M & 13 & S3 \\
US-37 & Exposer les facteurs aggravants et protecteurs de la prédiction & M & 8 & S3 \\
US-38 & Saisir les données oxymétriques et les index d'éveil & M & 5 & S4 \\
US-39 & Lancer une prédiction sans fichier, à partir d'une saisie manuelle & S & 8 & S4 \\
US-40 & Importer un rapport externe pour préremplir le formulaire clinique
& S & 5 & S4 \\
US-41 & Lire une interprétation clinique automatique des anomalies détectées
& S & 5 & S4 \\

\epictitre{E6 --- Visualisation, annotation et normes}
US-42 & Visualiser l'hypnogramme de façon interactive & M & 8 & S4 \\
US-43 & Voir la saturation et la densité d'événements sur le même axe temporel
& S & 5 & S4 \\
US-44 & Suivre la progression de l'analyse en cours & C & 5 & S4 \\
US-45 & Annoter une époque précise pour consigner une observation & S & 5 & S6 \\
US-46 & Situer chaque métrique par rapport aux normes de référence & S & 5 & S6 \\

\epictitre{E7 --- Collaboration entre praticiens}
US-47 & Ouvrir une discussion sur un examen avec un confrère & S & 8 & S6 \\
US-48 & Échanger des messages dans un fil rattaché à l'examen & S & 5 & S6 \\
US-49 & Consulter une boîte de réception centralisée & S & 5 & S6 \\

\epictitre{E8 --- Archivage et exports}
US-50 & Archiver le fichier brut sans bloquer l'interface & M & 8 & S5 \\
US-51 & Archiver l'hypnogramme et le compte rendu clinique & S & 5 & S5 \\
US-52 & Délivrer des liens d'accès à durée limitée & M & 5 & S5 \\
US-53 & Exporter les marqueurs calculés pour un autre outil & S & 5 & S4 \\
US-54 & Composer le compte rendu en choisissant les sections à inclure & M & 8 & S6 \\
US-55 & Synchroniser l'hypnogramme annoté vers le dossier & S & 5 & S6 \\

\epictitre{E9 --- Industrialisation et confort d'usage}
US-56 & Conteneuriser les trois services pour un déploiement reproductible
& M & 8 & S6 \\
US-57 & Orchestrer la pile en une seule commande & M & 5 & S6 \\
US-58 & Automatiser tests et construction d'images à chaque livraison & S & 8 & S6 \\
US-59 & Basculer l'interface entre français et anglais & S & 5 & S6 \\
US-60 & Basculer en thème sombre pour la lecture nocturne & C & 3 & S6 \\
\end{longtable}

\noindent
La formulation complète a été abrégée dans le tableau pour tenir sur une ligne ; le
bénéficiaire se déduit de l'épic.

\section{Diagramme de classes}

Le modèle métier s'organise autour de huit entités (figure~\ref{fig:classes}).
L'entité centrale est l'\emph{examen}, qui matérialise une nuit d'enregistrement : elle
porte les références vers les fichiers archivés, la classe de sévérité prédite et le
document de résultats. Autour d'elle gravitent le patient qui l'a subie, les annotations
que le praticien y a portées et les discussions qu'elle a suscitées.

\begin{figure}[H]
\centering
\begin{tikzpicture}[
    x=1cm,y=1cm,
    cls/.style  = {rectangle,draw=hypnoblue!65,fill=white,align=left,anchor=north,
                   font=\tiny,inner sep=0pt,text width=3.5cm},
    ent/.style  = {rectangle,draw=hypnored!70,fill=hypnored!6,align=left,anchor=north,
                   font=\tiny,inner sep=0pt,text width=3.5cm},
    rel/.style  = {hypnogray!85,line width=0.6pt},
    card/.style = {font=\tiny,text=hypnogray,inner sep=1.6pt,fill=white},
    note/.style = {font=\tiny\itshape,text=hypnogray,align=center,text width=3.5cm},
  ]

  \newcommand{\classe}[3]{%
    \begin{tabular}{@{\hspace{3pt}}p{3.3cm}@{}}
      \rowcolor{hypnolight}\textbf{#1}\\[1pt]
      \hline
      \rule{0pt}{7pt}#2\\[1pt]
      \hline
      \rule{0pt}{7pt}\textit{#3}\\[1pt]
    \end{tabular}}

  % ── Rangée 1 ────────────────────────────────────────────────────────────
  \node[cls] (hop) at (2.0, 4.40) {\classe{Hospital}{id, name, billing\_tier,
       b2\_bucket, created\_at}{--}};
  \node[cls] (usr) at (7.0, 4.40) {\classe{User}{id, email, hashed\_password, role,
       first\_name, last\_name, status, license\_expiry, last\_login}{login(),
       logout()}};
  \node[cls] (pat) at (12.0, 4.40) {\classe{Patient}{id, first\_name, last\_name,
       age, imc, gender}{--}};

  % ── Rangée 2 ────────────────────────────────────────────────────────────
  \node[cls] (log)  at (2.0, 0.55) {\classe{AuditLog}{id, timestamp, user\_email,
       user\_role, action, ip\_address}{--}};
  \node[cls] (conv) at (7.0, 0.55) {\classe{FileConversation}{id, file\_type,
       created\_at}{--}};
  \node[ent] (psg)  at (12.0, 0.55) {\classe{PSG}{id, date, severity, report\_data,
       edf\_url, hypnogram\_url, hypnogram\_annotated\_url, csv\_url,
       osa\_report\_url}{--}};

  % ── Rangée 3 ────────────────────────────────────────────────────────────
  \node[cls] (msg) at (7.0, -3.35) {\classe{FileMessage}{id, content, timestamp}{--}};
  \node[cls] (ann) at (12.0,-3.35) {\classe{PSGAnnotation}{id, epoch\_index, note,
       created\_at}{--}};

  \node[note,anchor=north] at (2.0,-1.05)
       {table volontairement dénormalisée :\\aucune clé étrangère, afin que la
        trace survive à la suppression du compte};

  % ── Associations ────────────────────────────────────────────────────────
  \draw[rel] (hop.east) -- node[card,pos=0.12] {1}
                           node[card,pos=0.88] {0..*} (usr.west);
  \draw[rel] (usr.east) -- node[card,pos=0.12] {1}
                           node[card,pos=0.88] {0..*} (pat.west);

  \draw[rel] (pat.south) -- node[card,pos=0.13] {1}
                            node[card,pos=0.87] {0..*} (psg.north);
  \draw[rel] (psg.south) -- node[card,pos=0.13] {1}
                            node[card,pos=0.87] {0..*} (ann.north);

  \draw[rel] (psg.west) -- node[card,pos=0.13] {1}
                           node[card,pos=0.87] {0..*} (conv.east);
  \draw[rel] (conv.south) -- node[card,pos=0.13] {1}
                             node[card,pos=0.87] {0..*} (msg.north);

  \draw[rel] (usr.south) -- node[card,pos=0.13] {2}
                            node[card,pos=0.88] {0..*} (conv.north);

  % Auteur d'un message : contournement par la gauche
  \draw[rel] (usr.west) -- ++(-0.55,0)
             node[card,pos=0.55] {1}
             |- node[card,pos=0.93] {0..*} (msg.west);

\end{tikzpicture}
\vspace{0.2cm}
\caption[Diagramme de classes du modèle métier]%
        {Modèle métier : l'examen est l'entité pivot du système.}
\label{fig:classes}
\end{figure}

Le tableau~\ref{tab:dico} décrit le rôle de chaque table. Le dictionnaire complet,
colonne par colonne, figure en annexe.

\begin{table}[H]
\centering

\small
\begin{tabular}{@{}lcL{8.5cm}@{}}
\toprule
\textbf{Table} & \textbf{Clés étrangères} & \textbf{Rôle et particularités} \\
\midrule
\texttt{hospitals} & --- &
Établissements de rattachement ; prévoit un compartiment de stockage propre à chaque
structure \\
\addlinespace[2pt]
\texttt{users} & \texttt{hospital\_id} &
Comptes unifiés pour les deux rôles ; porte le statut et l'échéance de licence,
tous deux évalués à la connexion \\
\addlinespace[2pt]
\texttt{patients} & \texttt{doctor\_id} &
Rattachement direct au praticien, qui fonde le cloisonnement des données \\
\addlinespace[2pt]
\texttt{psgs} & \texttt{patient\_id} &
Entité pivot ; ne stocke que des références vers les fichiers, jamais leur contenu \\
\addlinespace[2pt]
\texttt{psg\_annotations} & \texttt{psg\_id} &
Observations horodatées à l'époque ; suppression en cascade avec l'examen \\
\addlinespace[2pt]
\texttt{audit\_logs} & --- &
Journal volontairement dénormalisé : l'adresse électronique et le rôle sont recopiés,
afin que la trace survive à la suppression du compte \\
\addlinespace[2pt]
\texttt{file\_conversations} & \texttt{psg\_id}, deux \texttt{user\_id} &
Fil de discussion rattaché à un examen et à exactement deux praticiens \\
\addlinespace[2pt]
\texttt{file\_messages} & \texttt{conversation\_id}, \texttt{sender\_id} &
Messages horodatés ; suppression en cascade avec la discussion \\
\bottomrule
\end{tabular}
\caption{Tables du schéma relationnel}
\label{tab:dico}
\end{table}

Deux choix méritent un commentaire. Le \textbf{journal d'audit est dénormalisé} : il
recopie l'adresse et le rôle de l'auteur au lieu de référencer son compte. C'est
délibéré --- une trace de conformité doit rester lisible même après la suppression du
compte concerné, et une clé étrangère l'en empêcherait. Les \textbf{références de
fichiers}, quant à elles, ne sont jamais transmises telles quelles : elles sont
converties en liens signés à durée limitée au moment de la sérialisation de la réponse.

\section{Environnement de travail}

\subsection{Frameworks et langages utilisés}

Le tableau~\ref{tab:techno} récapitule les technologies retenues et la raison de leur
choix.

\begin{table}[H]
\centering\small
\begin{tabular}{@{}lL{4.0cm}L{6.6cm}@{}}
\hline
\ligneentete \entete{Couche} & \entete{Technologie} & \entete{Motif du choix} \\
\hline
Interface utilisateur & React 19 et Vite &
Composition par composants et rechargement instantané pendant le développement \\
Visualisation & Recharts &
Tracés interactifs superposables sur un axe temporel commun \\
Multilinguisme & react-i18next &
Bascule immédiate sans rechargement de l'application \\
Interface applicative & FastAPI &
Validation déclarative des données et documentation OpenAPI générée automatiquement \\
Apprentissage profond & PyTorch &
Écriture naturelle des architectures et accélération graphique \\
Apprentissage classique & scikit-learn, XGBoost, LightGBM &
Implémentations de référence et interface homogène pour la comparaison \\
Explicabilité & SHAP &
Attributions exactes en temps polynomial sur les modèles à arbres \\
Traitement du signal & MNE-Python, SciPy &
Lecture des formats physiologiques standards et rééchantillonnage \\
Persistance & PostgreSQL et SQLAlchemy &
Intégrité référentielle et correspondance objet-relationnel \\
Stockage des fichiers & Service objet compatible S3 &
Fichiers volumineux hors base, liens signés à durée limitée \\
Conteneurisation & Docker et Docker Compose &
Déploiement reproductible en une commande \\
\hline
\end{tabular}
\caption{Technologies retenues et motifs de choix}
\label{tab:techno}
\end{table}

\subsection{Environnement logiciel}

\begin{table}[H]
\centering

\small
\begin{tabular}{@{}lL{9.6cm}@{}}
\toprule
\textbf{Domaine} & \textbf{Outils} \\
\midrule
Langage & Python 3.10 \\
Traitement du signal & MNE-Python~\cite{gramfort2013}, SciPy, NumPy \\
Apprentissage profond & PyTorch~\cite{paszke2019} avec accélération CUDA \\
Apprentissage classique & scikit-learn~\cite{pedregosa2011},
XGBoost~\cite{chen2016}, LightGBM~\cite{ke2017} \\
Explicabilité & SHAP~\cite{lundberg2017,lundberg2020} \\
Expérimentation & Jupyter, matplotlib, seaborn \\
\bottomrule
\end{tabular}
\caption{Environnement logiciel}
\label{tab:environnement}
\end{table}

Le suivi de version s'appuie sur Git, la modélisation sur des diagrammes produits
directement dans la chaîne de composition du mémoire, et l'expérimentation sur des
carnets Jupyter conservés dans le dépôt à titre de trace reproductible.

\subsection{Environnement matériel}

Les entraînements ont été menés sur une station équipée d'un processeur graphique
NVIDIA GeForce RTX 4060 pour portable, disposant de \SI{8,6}{\giga\octet} de mémoire
vidéo, et de \SI{32}{\giga\octet} de mémoire vive. Une seconde station, équipée d'une
RTX 3060, a servi à la construction des ensembles.

Ce dimensionnement n'est pas anecdotique : il a imposé un arbitrage sur la quantité de
données exploitable, développé au chapitre~\ref{chap:sprint2}.

\section{Architecture globale}

\subsection{Architecture logique}

\hypnoria{} suit une architecture client-serveur à trois couches
(figure~\ref{fig:archi3t}). La couche de présentation est une application monopage
exécutée dans le navigateur du praticien. La couche applicative expose une interface de
programmation sans état, qui porte à la fois la logique métier et le moteur
d'inférence. La couche de persistance sépare délibérément deux natures de données : les
métadonnées cliniques, structurées et interrogeables, vont dans une base relationnelle ;
les fichiers volumineux partent vers un stockage objet.

\begin{figure}[H]
\centering
\begin{tikzpicture}[
    x=1cm,y=1cm,
    couche/.style = {rectangle,rounded corners=3pt,draw=hypnoblue!45,
                     fill=hypnolight!60,minimum width=13.6cm},
    mod/.style    = {rectangle,rounded corners=2pt,draw=hypnoblue!65,fill=white,
                     align=center,font=\scriptsize,minimum height=0.95cm,
                     text width=2.45cm},
    modml/.style  = {rectangle,rounded corners=2pt,draw=hypnored!70,fill=hypnored!7,
                     align=center,font=\scriptsize,minimum height=0.95cm,
                     text width=2.45cm},
    ext/.style    = {rectangle,rounded corners=2pt,draw=hypnogray!60,fill=white,
                     align=center,font=\scriptsize,minimum height=0.95cm,
                     text width=2.9cm},
    titre/.style  = {font=\scriptsize\bfseries,text=hypnoblue,anchor=west},
    fl/.style     = {{Latex[length=2mm]}-{Latex[length=2mm]},hypnoblue!70,
                     line width=0.8pt},
    lab/.style    = {font=\tiny\itshape,text=hypnogray,fill=white,inner sep=1.5pt},
  ]

  % ── Couche présentation ─────────────────────────────────────────────────
  \node[couche,minimum height=1.5cm] (c1) at (7.0,4.4) {};
  \node[titre] at (0.25,5.02) {Présentation};
  \node[mod] at (3.3,4.25) {Assistant\\d'analyse};
  \node[mod] at (6.0,4.25) {Registre\\patients};
  \node[mod] at (8.7,4.25) {Console\\d'administration};
  \node[mod] at (11.4,4.25) {Consultations};

  % ── Couche application ──────────────────────────────────────────────────
  \node[couche,minimum height=2.7cm] (c2) at (7.0,1.35) {};
  \node[titre] at (0.25,2.45) {Application};
  \node[mod]   at (3.3,1.75) {Authentification\\et rôles};
  \node[mod]   at (6.0,1.75) {Gestion\\métier};
  \node[mod]   at (8.7,1.75) {Archivage};
  \node[mod]   at (11.4,1.75) {Journalisation};
  \node[modml] at (4.65,0.45) {Prétraitement\\du signal};
  \node[modml] at (7.35,0.45) {Classification\\des stades};
  \node[modml] at (10.05,0.45) {Sévérité\\et explication};
  \node[font=\tiny\itshape,text=hypnored,anchor=east] at (3.30,0.45) {moteur d'inférence};

  % ── Couche persistance ──────────────────────────────────────────────────
  \node[couche,minimum height=1.5cm] (c3) at (7.0,-1.75) {};
  \node[titre] at (0.25,-1.13) {Persistance};
  \node[ext] at (4.6,-1.9) {Base relationnelle\\\tiny métadonnées cliniques};
  \node[ext] at (9.4,-1.9) {Stockage objet\\\tiny fichiers volumineux};

  % ── Flux ────────────────────────────────────────────────────────────────
  \draw[fl] (7.0,3.62) -- node[lab] {interface applicative sur HTTP} (7.0,2.75);
  \draw[fl] (4.6,-0.02) -- node[lab] {requêtes} (4.6,-1.12);
  \draw[fl] (9.4,-0.02) -- node[lab] {liens signés} (9.4,-1.12);

\end{tikzpicture}
\vspace{0.3cm}
\caption[Architecture générale du système]%
        {Architecture en trois couches. Les modules encadrés en rouge constituent le
         moteur d'inférence.}
\label{fig:archi3t}
\end{figure}

Quatre décisions structurantes méritent d'être argumentées ; le tableau~\ref{tab:choix}
les récapitule avec l'alternative écartée et le motif du choix.

\begin{table}[H]
\centering

\small
\begin{tabular}{@{}L{3.2cm}L{5.2cm}L{5.6cm}@{}}
\toprule
\textbf{Décision} & \textbf{Alternative écartée} & \textbf{Motif} \\
\midrule
Application monopage découplée de l'interface applicative &
Application web à rendu serveur &
L'analyse d'un enregistrement dure plusieurs dizaines de secondes ; un client
autonome peut afficher la progression, poursuivre les transferts en arrière-plan et
recomposer les visualisations sans rechargement \\
\addlinespace[3pt]
Moteur d'inférence hébergé dans le même processus que l'interface applicative &
Service d'inférence séparé &
Les modèles pèsent plusieurs centaines de mégaoctets en mémoire ; les charger une
seule fois et les conserver en cache évite un coût de démarrage à chaque requête. La
frontière logique reste posée pour une séparation ultérieure \\
\addlinespace[3pt]
Séparation métadonnées / fichiers volumineux &
Stockage des fichiers en base &
Un enregistrement pèse plusieurs centaines de mégaoctets. Le conserver en base
dégraderait les sauvegardes et les temps de réponse, sans bénéfice transactionnel \\
\addlinespace[3pt]
Authentification par jeton signé, sans session serveur &
Session en mémoire ou en base &
L'absence d'état côté serveur autorise la réplication horizontale sans mécanisme de
partage de session \\
\bottomrule
\end{tabular}
\caption{Décisions d'architecture et leur justification}
\label{tab:choix}
\end{table}

\subsection{Architecture physique}

La figure~\ref{fig:deploiement} décrit la topologie d'exécution. Chaque service est
conteneurisé et l'ensemble démarre par une commande unique. Le service de stockage est
le seul élément externe à l'infrastructure.

\begin{figure}[H]
\centering
\begin{tikzpicture}[
    x=1cm,y=1cm,
    noeud/.style  = {rectangle,rounded corners=3pt,draw=hypnoblue!55,fill=white,
                     align=center,font=\scriptsize,minimum height=1.5cm,
                     text width=3.2cm},
    hote/.style   = {rectangle,rounded corners=3pt,draw=hypnogray!55,dashed,
                     fill=hypnolight!35},
    cloud/.style  = {rectangle,rounded corners=8pt,draw=hypnogray!65,fill=white,
                     align=center,font=\scriptsize,minimum height=1.5cm,
                     text width=3.0cm},
    fl/.style     = {-{Latex[length=2mm]},hypnoblue!75,line width=0.8pt},
    lab/.style    = {font=\tiny\itshape,text=hypnogray,fill=white,inner sep=1.5pt},
  ]

  \node[hote,minimum width=11.4cm,minimum height=2.5cm] at (5.7,0) {};
  \node[font=\scriptsize\bfseries,text=hypnogray,anchor=north west] at (0.15,1.15)
       {Hôte de déploiement --- orchestration de conteneurs};

  \node[noeud] (front) at (2.1,-0.2)  {\textbf{Serveur web}\\\tiny Nginx\\\tiny
                                        application compilée};
  \node[noeud] (back)  at (5.7,-0.2)  {\textbf{Interface applicative}\\\tiny Uvicorn
                                        et FastAPI\\\tiny modèles en mémoire};
  \node[noeud] (db)    at (9.3,-0.2)  {\textbf{Base de données}\\\tiny PostgreSQL\\
                                        \tiny volume persistant};

  \node[cloud] (b2)    at (13.6,-0.2) {\textbf{Stockage objet}\\\tiny service externe\\
                                        \tiny interface S3};

  \node[font=\scriptsize\bfseries,align=center,text width=2.2cm] (nav) at (2.1,2.35)
       {Navigateur\\du praticien};

  \draw[fl] (nav) -- node[lab] {HTTPS} (front);
  \draw[fl] (front) -- node[lab] {mandataire} (back);
  \draw[fl] (back) -- node[lab] {SQL} (db);
  \draw[fl] (back) -- node[lab] {S3} (b2);

\end{tikzpicture}
\vspace{0.3cm}
\caption[Vue de déploiement]%
        {Topologie d'exécution : trois conteneurs sur un hôte, un service de stockage
         externe.}
\label{fig:deploiement}
\end{figure}

\section{Sprint Backlog général}

\subsection{Découpage en six sprints}

Les vingt-six semaines du stage ont été découpées en six sprints de quatre semaines,
suivis de deux semaines consacrées à la rédaction et à la préparation de la soutenance.
Le tableau~\ref{tab:sprints} présente l'objectif et le livrable de chacun.

\begin{table}[H]
\centering

\small
\begin{tabular}{@{}L{1.5cm}L{2.5cm}L{5.0cm}L{4.4cm}c@{}}
\toprule
\textbf{Sprint} & \textbf{Période} & \textbf{Objectif} & \textbf{Livrable} &
\textbf{Pts} \\
\midrule
S1 & 26 jan.\newline 22 fév. &
Cadrer le projet, lancer la demande d'accès aux cohortes et prouver la faisabilité de
la chaîne sans attendre l'autorisation &
Chaîne de prétraitement opérationnelle et premiers modèles de faisabilité & 29 \\
\addlinespace[3pt]
S2 & 23 fév.\newline 22 mars &
Passer du jeu d'amorçage à la cohorte clinique et couvrir les trois familles
d'architectures &
Douze modèles de base entraînés et évalués & 47 \\
\addlinespace[3pt]
S3 & 23 mars\newline 19 avr. &
Faire coopérer les architectures, puis franchir le pas du staging vers le diagnostic
expliqué &
Quatre ensembles et le modèle de sévérité avec ses explications & 42 \\
\addlinespace[3pt]
S4 & 20 avr.\newline 17 mai &
Rendre les modèles consommables et donner au praticien un parcours guidé &
API documentée et parcours complet du dépôt au compte rendu & 80 \\
\addlinespace[3pt]
S5 & 18 mai\newline 14 juin &
Transformer l'outil d'analyse en plateforme gouvernée et traçable &
Dossier patient persistant, archivage cloud, espace d'administration & 77 \\
\addlinespace[3pt]
S6 & 15 juin\newline 12 juil. &
Couvrir le travail à plusieurs, la restitution documentaire et la reproductibilité &
Collaboration, compte rendu exportable, pile conteneurisée & 83 \\
\bottomrule
\end{tabular}
\caption{Découpage du projet en six sprints}
\label{tab:sprints}
\end{table}

\subsection{Vélocité prévisionnelle}

La vélocité mesure le nombre de points effectivement livrés par sprint. Sa progression,
représentée à la figure~\ref{fig:velocite}, ne décrit pas un plateau régulier, et c'est
précisément ce qui la rend crédible : les deux premiers sprints sont absorbés par la
montée en compétence sur le cadre d'apprentissage profond et par des temps
d'entraînement incompressibles.

\begin{figure}[H]
\centering
\begin{tikzpicture}[x=1cm,y=1cm]

  % Axe des ordonnées
  \draw[hypnogray!60,line width=0.6pt] (0,0) -- (0,4.6);
  \foreach \v/\y in {0/0, 20/1.0, 40/2.0, 60/3.0, 80/4.0} {
    \draw[hypnogray!60] (-0.1,\y) -- (0,\y);
    \node[left,font=\scriptsize,text=hypnogray] at (-0.12,\y) {\v};
    \draw[hypnogray!18,line width=0.4pt] (0,\y) -- (10.6,\y);
  }
  \node[rotate=90,font=\scriptsize,text=hypnogray] at (-1.0,2.3) {points livrés};

  % Barres
  \foreach \i/\lab/\pts in {0/S1/29, 1/S2/47, 2/S3/42, 3/S4/80, 4/S5/77, 5/S6/83} {
    \pgfmathsetmacro{\xa}{0.55 + \i*1.65}
    \pgfmathsetmacro{\xb}{\xa + 1.1}
    \pgfmathsetmacro{\hh}{\pts/20}
    \fill[hypnoblue!65] (\xa,0) rectangle (\xb,\hh);
    \draw[hypnoblue] (\xa,0) rectangle (\xb,\hh);
    \node[font=\scriptsize\bfseries,text=hypnoblue,anchor=south]
         at ({(\xa+\xb)/2},\hh+0.04) {\pts};
    \node[font=\scriptsize,text=hypnogray,anchor=north]
         at ({(\xa+\xb)/2},-0.06) {\lab};
  }

  % Axe des abscisses
  \draw[hypnogray!60,line width=0.6pt] (0,0) -- (10.6,0);

  % Vélocité moyenne
  \draw[hypnored,dashed,line width=0.8pt] (0,2.983) -- (10.6,2.983);
  \node[font=\scriptsize\itshape,text=hypnored,anchor=west] at (9.0,3.28)
       {moyenne : 60};

\end{tikzpicture}
\vspace{0.3cm}
\caption[Vélocité par sprint]%
        {Points de complexité livrés par sprint.}
\label{fig:velocite}
\end{figure}

\subsection{Planning}

La figure~\ref{fig:gantt} confronte la planification aux travaux réellement menés. Le
jalon le plus structurant est l'obtention de l'accès aux cohortes en semaine 5. Le
recouvrement volontaire entre la demande d'accès et les travaux d'amorçage sur un jeu
public traduit une décision assumée : ne pas subir un délai administratif, mais
l'occuper.

\begin{figure}[H]
\centering
\begin{ganttchart}[
    x unit                  = 0.375cm,
    y unit title            = 0.52cm,
    y unit chart            = 0.44cm,
    hgrid                   = {draw=hypnogray!18},
    vgrid                   = {*1{draw=hypnogray!18}},
    title/.append style     = {draw=hypnoblue!45, fill=hypnolight},
    title label font        = \tiny,
    title height            = 1,
    bar/.append style       = {draw=hypnoblue!85, fill=hypnoblue!55,
                               rounded corners=1pt},
    bar height              = 0.58,
    bar label font          = \scriptsize,
    bar label node/.append style = {text=hypnoblue!95},
    milestone/.append style = {fill=hypnored, draw=hypnored!85},
    milestone label font    = \scriptsize\itshape,
    milestone label node/.append style = {text=hypnored},
    milestone height        = 0.42,
  ]{1}{26}

  \gantttitle{Sprint 1}{4}
  \gantttitle{Sprint 2}{4}
  \gantttitle{Sprint 3}{4}
  \gantttitle{Sprint 4}{4}
  \gantttitle{Sprint 5}{4}
  \gantttitle{Sprint 6}{4}
  \gantttitle{Clôture}{2} \\
  \gantttitlelist{1,...,26}{1} \\

  \ganttbar{Étude bibliographique}{1}{3} \\
  \ganttbar{Demande d'accès NSRR}{1}{5} \\
  \ganttbar{Chaîne de prétraitement}{2}{4} \\
  \ganttbar{Amorçage sur jeu public}{3}{4} \\
  \ganttmilestone{Accès aux cohortes obtenu}{5} \\
  \ganttbar{Préparation de la cohorte}{5}{6} \\
  \ganttbar{Modèles de base}{6}{9} \\
  \ganttmilestone{12 modèles entraînés}{9} \\
  \ganttbar{Ensembles par empilement}{9}{10} \\
  \ganttbar{Modèle de sévérité}{10}{12} \\
  \ganttbar{Service d'inférence}{13}{15} \\
  \ganttbar{Interface d'analyse}{14}{17} \\
  \ganttbar{Base de données, archivage}{17}{19} \\
  \ganttbar{Administration, audit}{19}{20} \\
  \ganttbar{Collaboration, annotation}{21}{22} \\
  \ganttbar{Compte rendu, multilinguisme}{22}{23} \\
  \ganttmilestone{Plateforme complète}{22} \\
  \ganttbar{Conteneurisation, intégration}{23}{24} \\
  \ganttbar{Rédaction du mémoire}{21}{26} \\
  \ganttmilestone{Soutenance}{26}

\end{ganttchart}
\vspace{0.3cm}
\caption[Diagramme de Gantt du projet]%
        {Déroulement du projet sur les vingt-six semaines du stage.}
\label{fig:gantt}
\end{figure}

\sectionsansnum{Conclusion}

Ce chapitre a traduit les intentions du chapitre précédent en une commande vérifiable.
Quatre acteurs ont été identifiés, les besoins fonctionnels regroupés en neuf épics et
déclinés en soixante récits utilisateur, et les besoins non fonctionnels énoncés avec
leur moyen de vérification. Le modèle de données, l'environnement de travail et
l'architecture globale ont été présentés.

L'ordonnancement du backlog a produit six sprints de quatre semaines. Les six chapitres
suivants leur sont consacrés, chacun suivant la même trame : objectif du sprint, analyse
--- backlog, cas d'utilisation, diagrammes de séquence --- puis réalisation et bilan.
